\section*{Capítulo \ref{ch:g12:ph-conductivity-titrations} --- Valoraciones por pH y por conductividad}

\begin{solution}{exo:g12:ph-conductivity-titrations:1}
$V_E = \qty{20.0}{mL}$; $c = 0.100 \times 20.0 / 20.0 = \qty{0.100}{mol/L}$.
\end{solution}

\begin{solution}{exo:g12:ph-conductivity-titrations:2}
En la \omterm{def:g12:ph-conductivity-titrations:ph-metric}{semiequivalencia}, \omterm{def:g9:acids-bases-ph:ph}{pH} $= pK_a$: $pK_a = 3.75$, próximo al del ácido
metanoico (3.74).
\end{solution}

\begin{solution}{exo:g12:ph-conductivity-titrations:3}
La fenolftaleína: su \omterm{def:g12:buffers-predominance:indicator}{zona de viraje}, 8.0--10.0, está dentro del salto,
que es básico. El rojo de metilo (4.2--6.3) viraría cerca de la
\omterm{def:g12:ph-conductivity-titrations:ph-metric}{semiequivalencia}, demasiado pronto.
\end{solution}

\begin{solution}{exo:g12:ph-conductivity-titrations:4}
La respuesta de la sonda deriva con el tiempo y la temperatura: se ajusta
para que lea correctamente en dos \omterm{def:g12:buffers-predominance:buffer}{disoluciones tampón} de \omterm{def:g9:acids-bases-ph:ph}{pH} conocido.
\end{solution}

\begin{solution}{exo:g12:ph-conductivity-titrations:5}
\ce{H3O+} (349.8) > \ce{OH-} (198.3) > \ce{Cl-} (76.2) > \ce{Na+} (50.3).
\end{solution}

\begin{solution}{exo:g12:ph-conductivity-titrations:6}
$c = 0.050 \times 14.6 / 20.0 = \qty{0.0365}{mol/L}$.
\end{solution}

\begin{solution}{exo:g12:ph-conductivity-titrations:7}
Pendientes: 2.0, 4.5, 17.5 y \qty{2.5}{mL^{-1}}. La mayor está entre
10.0 y \qty{10.2}{mL}: $V_E \approx \qty{10.1}{mL}$.
\end{solution}

\begin{solution}{exo:g12:ph-conductivity-titrations:8}
Antes de la \omterm{def:g11:titration:equivalence}{equivalencia}, cada \ce{OH-} añadido elimina un \ce{H3O+} (muy
conductor) y aporta un \ce{Na+} (mucho menos conductor): la \omterm{def:g12:ph-conductivity-titrations:conductivity}{conductividad}
baja. Después, el \ce{Na+} y el \ce{OH-} se acumulan sin reaccionar:
sube.
\end{solution}

\begin{solution}{exo:g12:ph-conductivity-titrations:9}
Al principio, pocos \omterm{def:g9:ions:ion}{iones} (el \omterm{def:g12:ka-and-pka:strong-acid}{ácido débil} apenas reacciona con el agua):
\omterm{def:g12:ph-conductivity-titrations:conductivity}{conductividad} baja. Antes de la \omterm{def:g11:titration:equivalence}{equivalencia} se forman \omterm{def:g9:ions:ion}{iones} etanoato y
sodio en lugar de \omterm{def:g7:atoms-and-molecules:molecule}{moléculas} sin carga: sube, despacio. Después de la
\omterm{def:g11:titration:equivalence}{equivalencia}, el \ce{Na+} y el \ce{OH-} se acumulan: sube más deprisa.
Dos rectas crecientes, la segunda más empinada.
\end{solution}

\begin{solution}{exo:g12:ph-conductivity-titrations:10}
Alrededor de 2.9, 4.76 y 8.7. El ácido etanoico es débil: solo un pequeño
porcentaje da \omterm{def:g12:ka-and-pka:oxonium}{iones oxonio}, así que el \omterm{def:g9:acids-bases-ph:ph}{pH} empieza más alto que el 1.0 del
ácido clorhídrico a una concentración parecida.
\end{solution}

\begin{solution}{exo:g12:ph-conductivity-titrations:11}
Para que el volumen de \omterm{def:g11:titration:titration}{valorante} añadido apenas cambie el volumen total:
así las concentraciones, y por tanto la \omterm{def:g12:ph-conductivity-titrations:conductivity}{conductividad}, varían siguiendo
rectas.
\end{solution}

\begin{solution}{exo:g12:ph-conductivity-titrations:12}
Los \omterm{def:g9:ions:ion}{iones} hidróxido reaccionan primero con el \omterm{def:g12:ka-and-pka:strong-acid}{ácido fuerte} (\ce{H3O+}) y
después con el débil. El primer salto marca el final del ácido
clorhídrico (su \omterm{def:g11:titration:equivalence}{volumen de equivalencia} lo mide), y el segundo, el final
del ácido etanoico (el volumen entre los dos saltos lo mide).
\end{solution}

\begin{solution}{exo:g12:ph-conductivity-titrations:13}
Al principio: $(349.8 + 76.2) \times 0.0100 = \qty{4.26}{mS/cm}$. En la
\omterm{def:g11:titration:equivalence}{equivalencia}: $[\ce{Na+}] = [\ce{Cl-}] = 1.00/110 = \qty{0.00909}{mol/L}$,
así que $(50.3 + 76.2) \times 0.00909 = \qty{1.15}{mS/cm}$. Ambos valores
coinciden con la figura.
\end{solution}

\begin{solution}{exo:g12:ph-conductivity-titrations:14}
Ninguno en $V_E$: el salto está en el mismo volumen, sea cual sea el
desfase de las lecturas. El $pK_a$ leído en la \omterm{def:g12:ph-conductivity-titrations:ph-metric}{semiequivalencia} sale 0.2
demasiado alto.
\end{solution}

\begin{solution}{exo:g12:ph-conductivity-titrations:15}
El \omterm{def:g9:acids-bases-ph:ph}{pH} no cambia durante esta \omterm{def:g9:ions:precipitate}{precipitación}, pero los \omterm{def:g9:ions:ion}{iones} sí: antes de
la \omterm{def:g11:titration:equivalence}{equivalencia}, cada \ce{Ag+} añadido elimina un \ce{Cl-} y aporta un
\ce{NO3-} de \omterm{def:g12:ph-conductivity-titrations:conductivity}{conductividad} parecida, así que la \omterm{def:g12:ph-conductivity-titrations:conductivity}{conductividad} se mantiene
casi constante (los \omterm{def:g9:ions:ion}{iones} sodio estaban desde el principio); después de
la \omterm{def:g11:titration:equivalence}{equivalencia}, el \ce{Ag+} y el \ce{NO3-} se acumulan y sube. Una recta
horizontal y después una creciente.
\end{solution}

\begin{solution}{pb:g12:ph-conductivity-titrations:1}
\textbf{1.} \qty{6}{g} de ácido etanoico por cada \qty{100}{g} de vinagre.

\textbf{2.} \qty{101}{g} de vinagre contienen $6.0 \times 101/100 =
\qty{6.06}{g}$; $M = \qty{60.0}{g/mol}$: \qty{0.101}{mol} en
\qty{0.1000}{L}, es decir, \qty{1.01}{mol/L}.

\textbf{3.} Sin diluir, \qty{10.0}{mL} necesitarían unos \qty{101}{mL} de
hidróxido de sodio a \qty{0.100}{mol/L}: más de lo que cabe en una
bureta.

\textbf{4.} Una pipeta aforada de \qty{10.0}{mL} y un matraz aforado de
\qty{100.0}{mL}.

\textbf{5.} \ce{CH3COOH + OH- -> CH3COO- + H2O}.

\textbf{6.} $V_E = \qty{10.1}{mL}$.

\textbf{7.} $c = 0.100 \times 10.1 / 10.0 = \qty{0.101}{mol/L}$.

\textbf{8.} $10 \times 0.101 = \qty{1.01}{mol/L}$.

\textbf{9.} \omterm{def:g9:acids-bases-ph:ph}{pH} 4.76 a \qty{5.05}{mL}: el $pK_a$ del ácido etanoico.

\textbf{10.} En la \omterm{def:g11:titration:equivalence}{equivalencia}, la disolución contiene etanoato de sodio,
y el etanoato es una \omterm{def:g12:ka-and-pka:strong-acid}{base débil}.

\textbf{11.} La fenolftaleína, cuya \omterm{def:g12:buffers-predominance:indicator}{zona de viraje}, 8.0--10.0, contiene el
\omterm{def:g9:acids-bases-ph:ph}{pH} en la \omterm{def:g11:titration:equivalence}{equivalencia} (alrededor de 8.7).

\textbf{12.} \omterm{def:g9:ions:ion}{Iones} sodio e \omterm{def:g9:ions:ion}{iones} etanoato.

\textbf{13.} Cada \ce{OH-} añadido convierte una \omterm{def:g7:atoms-and-molecules:molecule}{molécula} casi sin
ionizar en \ce{CH3COO-}, acompañado de un \ce{Na+}: se añaden lentamente
\omterm{def:g9:ions:ion}{iones} de \omterm{def:g12:ph-conductivity-titrations:conductivity}{conductividad} moderada.

\textbf{14.} Pasada la \omterm{def:g11:titration:equivalence}{equivalencia}, los \omterm{def:g9:ions:ion}{iones} hidróxido, muy
conductores, se acumulan junto con los \omterm{def:g9:ions:ion}{iones} sodio.

\textbf{15.} El \omterm{def:g12:ka-and-pka:strong-acid}{ácido débil} da muy pocos \omterm{def:g9:ions:ion}{iones} al agua.

\textbf{16.} En $V_E = \qty{10.1}{mL}$, como en la \omterm{def:g12:ph-conductivity-titrations:ph-metric}{valoración pH-métrica}.

\textbf{17.} $1.01 \times 60.0 = \qty{60.6}{g/L}$.

\textbf{18.} \qty{6.06}{g} en \qty{100.0}{mL}.

\textbf{19.} $6.06 \times 100 / 101 = \qty{6.0}{g}$ por cada
\qty{100}{g}.

\textbf{20.} El vinagre tiene una acidez de \qty{6.0}{\percent}: la etiqueta es
correcta.
\end{solution}
